\chapter{Infraestrutura}

O \textit{campus} Londrina da UTFPR dispõe de infraestrutura física e tecnológica compatível com as demandas dos cursos de graduação e pós-graduação ofertados, atendendo atualmente aproximadamente 1.906 estudantes regularmente matriculados, com o suporte de 163 docentes e 68 servidores técnico-administrativos. Inserido em uma área total de 109.561,46 m², com 24.174,20 m² de área construída, o campus oferece condições estruturais adequadas ao desenvolvimento das atividades acadêmicas e administrativas. No que se refere ao curso de CDIA, a infraestrutura disponível está alinhada aos objetivos formativos, ao perfil do egresso e às competências profissionais estabelecidas no Projeto Pedagógico de Curso (PPC), garantindo suporte integrado às atividades de ensino, pesquisa e extensão, bem como às demandas do corpo docente, discente e da coordenação. Esta seção apresenta a infraestrutura de apoio acadêmico disponibilizada à comunidade do curso.

\section{Espaços de Trabalho e Desenvolvimento de Atividades Docente e Coordenador de Curso}

O curso dispõe de espaços de trabalho destinados à coordenação, situada na sala K101-8, e ao corpo docente, alocado nas salas K101-8 e K101-9 do campus Londrina. Esses ambientes são ocupados de forma compartilhada e adequados às necessidades do curso, considerando o número de docentes, a carga horária acadêmica e as especificidades da área, viabilizando as ações acadêmico-administrativas, pedagógicas e de gestão.

Os ambientes contam com mobiliário, equipamentos de informática, acesso à rede institucional, ao Sistema Acadêmico da UTFPR, ao e-mail institucional e ao Ambiente Virtual de Aprendizagem Moodle e recursos de webconferência. A infraestrutura tecnológica disponível possibilita diferentes formas de organização do trabalho docente, incluindo atividades presenciais e remotas, tais como planejamento de ensino, orientação acadêmica, pesquisa e extensão. O acesso remoto aos sistemas institucionais é realizado por meio das plataformas oficiais da universidade, sendo que a instituição disponibiliza equipamentos para trabalho remoto conforme regulamentação interna e responsabilidades definidas nos instrumentos normativos vigentes.

O coordenador do curso dispõe de ambiente que permite o atendimento individual ou em grupo com privacidade. As salas docentes também possibilitam atendimento a estudantes, preparação de aulas, orientações presenciais e realização de reuniões pedagógicas, fornecendo condições para o desenvolvimento da atividade pedagógica presencial.

Há estrutura técnico-administrativa centralizada no campus, incluindo a Secretaria de Gestão Acadêmica (SEGEA-LD e  DERAC-LD), responsável por registros acadêmicos, matrículas, organização de horários e emissão de documentos, e a Coordenadoria de Gestão de Tecnologia da Informação (COGETI-LD), responsável pela infraestrutura de TIC, suporte técnico e manutenção dos sistemas institucionais. Essa organização concede apoio técnico-administrativo estruturado e funcionamento das rotinas acadêmicas.

Os ambientes destinados às atividades docentes e de coordenação dispõem de locais para a guarda de materiais pessoais durante o período de permanência no campus. Para momentos de pausa e descanso em atividades presenciais, os docentes utilizam espaços institucionais comuns, incluindo copa, cantina, restaurante universitário e áreas externas estruturadas com arborização, bancos e áreas gramadas, favorecendo bem-estar e integração.

A infraestrutura física e tecnológica está vinculada a instrumentos institucionais de planejamento e gestão, que contemplam manutenção preventiva e corretiva periódica, bem como planejamento de melhorias e adequações. A universidade mantém contratos terceirizados para manutenção de infraestrutura predial, elétrica e especializada. As solicitações de manutenção são registradas em sistema institucional direcionado ao Departamento de Serviços Gerais (DESEG-LD), proporcionando controle formal, rastreabilidade e acompanhamento das demandas. As solicitações relacionadas à Tecnologia da Informação são igualmente registradas em sistema, permitindo priorização técnica e monitoramento até sua conclusão.

Os ambientes acadêmicos e administrativos observam princípios de acessibilidade, ergonomia e segurança atendendo a Lei no 13.146 de 06 de julho de 2015 que institui a Lei Brasileira de inclusão da pessoa com deficiência \cite{Brasil_2015_Deficiencia}. O campus dispõe de rampas, elevadores, sanitários adaptados, sinalização adequada e demais recursos arquitetônicos conforme normas técnicas vigentes. No âmbito digital, o Moodle institucional possui recursos de acessibilidade compatíveis com diretrizes internacionais, permitindo navegação por leitores de tela, ajustes de contraste e ampliação de fontes (WORLD WIDE WEB CONSORTIUM, 2023; BRASIL, 2023). O campus conta com o Núcleo de Acessibilidade e Inclusão (NAI-LD), responsável por articular ações de inclusão, adaptação pedagógica e eliminação de barreiras, oferecendo condições de acesso, participação e aprendizagem a estudantes com deficiência ou necessidades educacionais específicas.

A avaliação dos espaços físicos, dos serviços institucionais e das condições de trabalho integra o processo contínuo de avaliação institucional conduzido pela Comissão Própria de Avaliação (CPA) da UTFPR \cite{UTFPR_CPA_2004, UTFPR_CPA_2009}. Esse processo contempla instrumentos periódicos de coleta de feedback de docentes, discentes e técnicos-administrativos, incluindo avaliação docente, avaliação de desempenho e pesquisas de clima organizacional, cujos resultados subsidiam ações de melhoria contínua da infraestrutura e dos serviços.

O campus também propicia acesso a programas institucionais de qualidade de vida e bem-estar, desenvolvidos em consonância com as diretrizes da gestão de pessoas da UTFPR (UTFPR, 2018a). As ações são devidamente publicadas por meio de correio eletrônico institucional e murais físicos para ampla divulgação nos ambientes de trabalho.

Dessa forma, a infraestrutura disponibilizada ao curso atende às exigências relativas aos espaços de trabalho para docentes e coordenação, viabilizando suporte tecnológico, condições adequadas de atendimento, manutenção periódica, planejamento de melhorias, inclusão e acessibilidade, bem como apoio institucional necessário ao desenvolvimento das atividades acadêmicas presenciais e remotas.

\section{Equipamentos e Serviços para Docentes e Coordenador de Curso}

Os equipamentos e serviços disponibilizados ao curso promovem acesso a recursos digitais compatíveis com as diretrizes estabelecidas no PPC e apoiam às atividades de ensino, pesquisa, extensão e gestão acadêmica.

A infraestrutura tecnológica compreende rede cabeada e cobertura integral de rede sem fio institucional, acessível mediante autenticação com credenciais institucionais. A conectividade apresenta capacidade para suportar simultaneamente atividades acadêmicas, acesso ao Ambiente Virtual de Aprendizagem (Moodle), utilização de sistemas institucionais e realização de webconferências. O curso utiliza sistemas de webconferência por meio do Google Meet institucional e da plataforma da Rede Nacional de Ensino e Pesquisa (RNP), oferecendo condições adequadas de interação síncrona. O Moodle institucional encontra-se hospedado no datacenter institucional da UTFPR promovendo estabilidade operacional e controle sobre a infraestrutura.

As solicitações de suporte tecnológico são registradas em sistema institucional formal, que permite categorizar demandas (preventivas, corretivas, emergenciais e outras), acompanhamento do atendimento e rastreabilidade das intervenções, caracterizando mecanismo estruturado de gestão de serviços de TI. O modelo adotado permite priorização técnica e organização do fluxo de atendimento.

Os equipamentos e serviços possibilitam acesso a aplicativos e softwares especializados necessários ao desenvolvimento do curso, incluindo Python, R, MatLab, Jupyter, Anaconda, Visual Studio Code, TensorFlow, PyTorch, MySQL, PostgreSQL, GitHub e Docker, entre outros. Tais ferramentas encontram-se instaladas nos laboratórios didáticos e, em sua maioria, são de licença livre, possibilitando também utilização remota pelos estudantes e docentes. Esse conjunto de recursos certificam aderência às competências técnicas previstas no PPC.

O acesso a informações acadêmicas ocorre por meio do Sistema Acadêmico da UTFPR, atualizado periodicamente, permitindo consulta a matrícula, desempenho, registros acadêmicos e demais informações institucionais. O Ambiente Virtual de Aprendizagem Moodle\footnote{https://moodle.utfpr.edu.br} disponibiliza fóruns, mensagens internas, envio de atividades, avaliações, feedback digital e difusão de conteúdos multimídia, constituindo plataforma de apoio às atividades acadêmicas e à comunicação entre docentes e discentes.

No que se refere à segurança da informação e ao tratamento de dados, a UTFPR possui Política formal de Segurança da Informação e Comunicação, inserida em sua estrutura de governança de tecnologia \cite{UTFPR_Seguranca_2018}. A instituição encontra-se em processo contínuo de adequação às disposições da Lei Geral de Proteção de Dados Pessoais (LGPD), com revisão de processos, sistemas e capacitação institucional.

A universidade mantém convênio com o Google Workspace institucional\footnote{https://ajuda.utfpr.edu.br/pt-br/servicos\_deinfra/gsuite}, que disponibiliza serviços de armazenamento em nuvem com autenticação institucional única, controle de acesso e mecanismos de proteção de dados. Os serviços digitais disponibilizados pela instituição são acessíveis externamente às suas instalações, permitindo acesso remoto ao Sistema Acadêmico, ao Moodle, ao e-mail institucional, à biblioteca digital e demais plataformas oficiais. Essa característica garante compatibilidade com as necessidades pedagógicas e administrativas do curso, possibilitando continuidade das atividades acadêmicas em ambientes presenciais e remotos.

Logo, os equipamentos e serviços providos ao curso demonstram compatibilidade com as exigências pedagógicas e administrativas, concedem acesso a recursos digitais especializados, garantem suporte tecnológico estruturado e observam princípios de segurança da informação e acessibilidade.

\section{Salas de Aulas e Demais Ambientes de Ensino-Aprendizagem (Exceto Laboratórios)}

O curso utiliza regularmente oito salas de aula teóricas com capacidade de 44 lugares, todas equipadas com projetor multimídia fixo com tela retrátil, computador para o professor, carteiras e cadeiras individuais móveis, espaços adequados para cadeirantes, ar-condicionado, ventilação natural, iluminação predominante com lâmpadas LED, sinalização de emergência e equipamentos de segurança. Algumas salas dispõem de quadro com giz e outras de lousa branca.

A manutenção das salas segue o sistema institucional de gestão e chamados atendidos pelo Departamento de Serviços Gerais (DESEG) por meio de uma plataforma de apoio informatizada de solicitação de demandas de serviços\footnote{https://apoio.utfpr.edu.br}, garantindo a conservação dos equipamentos e do mobiliário. Além disso, a infraestrutura é avaliada periodicamente pela Comissão Própria de Avaliação (CPA), incluindo conforto térmico e acústico, iluminação, recursos multimídia, limpeza, conservação e acessibilidade, permitindo monitoramento contínuo da qualidade do ambiente de ensino.

As salas possuem capacidade compatível com o número de estudantes por turma, garantindo que cada aluno tenha uma carteira individual. Além disso, o planejamento acadêmico considera o tamanho das turmas na alocação das salas, de modo a otimizar o aproveitamento do espaço físico disponível. Todas as salas possuem acesso à rede sem fio funcional, permitindo que alunos e docentes utilizem seus próprios notebooks e demais equipamentos que necessitem conexão à Internet, com espaço adequado no mobiliário e tomadas suficientes para conexão de dispositivos.

Também é possível reorganizar o mobiliário para atividades em grupos, metodologias ativas, sala invertida, workshops e discussões. Existe também sala com carteiras modulares, adequada para trabalho colaborativo e eventos interativos. As salas oferecem projetores multimídia e equipamentos para videoconferência, permitindo a realização de apresentações, aulas síncronas e interação com participantes externos, inclusive em formato presencial-remoto. O uso institucional do \textit{Google Meet} possibilita projeção coletiva e participação de docentes e alunos de outros campi ou de instituições parceiras.

O curso faz uso de \textit{Google Workspace} para e-mails, armazenamento de arquivos e videoconferências, bem como do Moodle em tempo real, possibilitando o acesso a materiais digitais, aplicação de quizzes, submissão de atividades e exercícios de fixação. Ferramentas colaborativas como \textit{Google Docs}, \textit{Google Drive}, \textit{Google Meet}, \textit{GitHub Education} e \textit{Google Colab} são frequentemente utilizadas para produção de conteúdo, desenvolvimento e compartilhamento de códigos e experimentos em Ciência de Dados e Inteligência Artificial.

Além disso, é possível realizar o compartilhamento de tela em grupo, permitindo que múltiplos participantes exibam conteúdo de seus dispositivos em uma única tela, por meio de Google Meet ou a plataforma de Conferência Web da RNP. Esses recursos oferecem suporte a projetos colaborativos com outras instituições e eventos com participação remota, incluindo bancas avaliadoras com membros externos, disciplinas compartilhadas entre campi e aulas ministradas por professores convidados de outras instituições.

Assim, o curso provê infraestrutura adequada, conectividade e recursos tecnológicos que favorecem metodologias ativas, ensino colaborativo e integração com o ecossistema acadêmico nacional e internacional.

\subsection{Ambientes de Aprendizagem em Educação Profissional e Tecnológica}

O curso dispõe de ambientes de aprendizagem planejados para desenvolvimento de competências profissionais e tecnológicas, simulando cenários de aplicação real. Esses ambientes permitem que os estudantes integrem teoria e prática em projetos de análise de dados, construção de sistemas de recomendação, desenvolvimento de algoritmos de aprendizado de máquina, processamento de linguagem natural e visão computacional, promovendo habilidades técnicas, analíticas e de resolução de problemas. O Quadro \ref{quadro:4.1} apresenta os ambientes voltados para o ensino tecnológico das tecnologias voltadas para CD/IA.

\begin{quadro}[htb]
    \centering
    \caption{Ambientes de ensino e aprendizagem}
    \label{quadro:4.1}
    \vspace{0.2cm}
    \begin{tabular}{|l|c|c|c|}
        \hline
        \textbf{Laboratório} & \textbf{Local} & \textbf{Área ($m^2$)} & \textbf{Capacidade (alunos)} \\ \hline
        Laboratório de Informática 1 & K-106 & 78 & 40 \\ \hline
        Laboratório de Informática 2 & K-107 & 78 & 40 \\ \hline
        Laboratório de Informática 3 & K-109 & 78 & 40 \\ \hline
        Laboratório de Informática 4 & K-102 & 25 & 24 \\ \hline
        Laboratório de Informática 5 & T-106 & 78 & 40 \\ \hline
        Laboratório de CAD & K-113 & 78 & 40 \\ \hline
        Laboratório de Eletrotécnica & K-007 & 80 & 24 \\ \hline
    \end{tabular}
    %\fonte{Autoria prória (2026)}
    \small \\ Fonte: Autoria própria (2026)
\end{quadro}

Cada laboratório é equipado com computadores, GPUs, softwares especializados (Python, R, Jupyter, TensorFlow, PyTorch) e bases de dados reais ou simuladas, permitindo aos estudantes executar experimentos e projetos que reproduzem situações do mercado de trabalho. Exemplos de atividades realizadas nos ambientes:
\begin{itemize}
    \item \textbf{Laboratórios de Informática:} programação, análise de dados, construção de modelos preditivos e pipelines de IA;
    \item \textbf{Laboratório de CAD:} desenvolvimento de visualizações gráficas de dados e projetos de interface de sistemas inteligentes;
    \item \textbf{Laboratório de Eletrotécnica:} integração de hardware com sistemas inteligentes, Internet das Coisas e prototipagem de sensores e atuadores.
\end{itemize}

A utilização dos ambientes é planejada e supervisionada, com capacitação docente e manutenção periódica dos recursos, de forma que os estudantes possam desenvolver competências profissionais em contextos próximos à prática real.

\section{Acesso dos Estudantes a Equipamentos de Informática}

O curso dispõe de infraestrutura de informática adequada às necessidades do curso e dos estudantes e promove atividades práticas, experimentações e desenvolvimento de projetos colaborativos.

Todos os laboratórios do curso possuem computadores com configuração padronizada conforme apresentado no Quadro \ref{quadro:4.2}, equipados com softwares essenciais para a formação, incluindo navegador web, Visual Studio Code, Python, R, Matlab, ferramentas de gerenciamento de bancos de dados (PostgreSQL e MySQL), Jupyter e GitHub, entre outros. Os laboratórios estão adequados ao espaço físico disponível, permitindo que cada aluno utilize um computador individualmente. Todos os ambientes possuem rede sem fio, compatível com os padrões 2.4 GHz e 5 GHz, com velocidade de transmissão de aproximadamente 400 Mbps, oferecendo banda suficiente para atender às demandas acadêmicas.

\begin{quadro}[htb]
    \centering
    \caption{Configuração dos computadores do laboratório do curso CD/IA}
    \label{quadro:4.2}
    \vspace{0.2cm}
    \begin{tabular}{|c|c|c|c|}
        \hline
        \textbf{Laboratório} & \textbf{CPU} & \textbf{Memória RAM} & \textbf{Armazenamento} \\ \hline
        K-102 & Intel Core i7-14700 & 16GB & 1 TB \\ \hline
        K-106 & AMD A8-5500B & 8GB & 480 GB \\ \hline
        K-107 & Intel i7-7700 & 16GB & 512GB \\ \hline
        K-109 & Intel i5-9500 & 16GB & 512GB \\ \hline
        K-113 & Intel i5-4590 & 12GB & 480GB \\ \hline
        T-106 & AMD Ryzen 7 5700G & 16GB & 480GB \\ \hline
    \end{tabular}
    %\fonte{Autoria prória (2026)}
    \small \\ Fonte: Autoria própria (2026)
\end{quadro}

A infraestrutura de informática é avaliada periodicamente pelos estudantes, por meio de questionários formais conduzidos no âmbito da UTFPR e alinhados às 10 dimensões do SINAES, contemplando especificamente a qualidade, adequação e funcionamento dos laboratórios. O campus conta com a Coordenadoria de Gestão de Tecnologia da Informação (COGETI-LD), que disponibiliza equipe de suporte técnico responsável pela manutenção física e lógica dos computadores, instalação de softwares específicos, gerenciamento da rede institucional e suporte a equipamentos de TI.

O uso dos laboratórios é organizado de forma que cada estudante tenha acesso a um computador individual, possibilitando estudos individuais e em grupo, com horários de utilização restritos às aulas ou sob supervisão de docentes responsáveis. A infraestrutura atende à demanda do curso em termos de quantidade, atualização e adequação dos equipamentos, certificando qualidade de funcionamento e pertinência aos propósitos do curso. Os laboratórios suportam execução de experimentações práticas, desenvolvimento de algoritmos, análise de dados, modelagem e demais atividades previstas no PPC.

Acrescenta-se aos recursos tecnológicos, a formação continuada para docentes, técnicos e estudantes, incluindo cursos de extensão, monitorias, semanas acadêmicas, capacitações internas pelo AVA (Moodle), Escola Virtual do Governo e Capacita Gov.br. Embora a UTFPR não possua diretrizes institucionais específicas sobre o uso de tecnologias emergentes ou Inteligência Artificial no ensino, o curso adota boas práticas pedagógicas e incentiva a utilização de softwares e recursos tecnológicos atualizados, permitindo que os estudantes se familiarizar com ferramentas emergentes e desenvolvam competências práticas em CDIA.


\section{Tecnologias e Recursos Digitais}

O curso utiliza diversas plataformas digitais integradas ao ambiente acadêmico, incluindo AVA Moodle, \textit{Google Workspace}, \textit{Google Colab}, \textit{GitHub Education}, \textit{Microsoft Azure}, \textit{Office 365}, \textit{Overleaf} institucional, Certificado Digital Pessoal ICPEdu, Portal de Informação em Acesso Aberto (PIAA), Repositório Institucional da UTFPR (RIUT), Portal de Periódicos Científicos da UTFPR (PERI), Comunidade Acadêmico Federada (CAFe), entre outros. Essas ferramentas são complementadas pelo Sistema Acadêmico da UTFPR, que oferece controle completo da vida acadêmica, permitindo registro de notas, acompanhamento de disciplinas e gestão de processos administrativos e pedagógicos. Todos os recursos digitais são de amplo acesso, devidamente documentados e disponibilizados de forma clara aos usuários por meio de uma página institucional centralizada\footnote{https://utfpr.edu.br}.

Além disso, a UTFPR garante acessibilidade digital: o portal institucional conta com barra de acessibilidade conforme padrão e-MAG/WCAG, e o Moodle é utilizável por pessoas com deficiências visuais e auditivas, permitindo inclusão e usabilidade para todos os estudantes.

Esses recursos digitais são integrados às atividades acadêmicas do curso, possibilitando apresentações de TCC, projetos, seminários e workshops. Nas salas de aula, projetores multimídia e ferramentas de videoconferência, como Google Meet e Conferência Web da RNP, permitem aulas híbridas e apresentações síncronas. Dessa forma, as plataformas digitais favorecem a integração entre conteúdos teóricos e atividades práticas, proporcionando ensino contextualizado da aplicação de conceitos de Ciência de Dados e Inteligência Artificial.

Para promover ensino colaborativo, o AVA Moodle fornece fóruns, chats e repositórios de arquivos assíncronos, integrados às plataformas de videoconferência (\textit{Google Meet}, \textit{Microsoft Teams}, RNP). Suítes de produtividade na nuvem, como \textit{Microsoft 365} e \textit{Google Workspace}, permitem compartilhamento de arquivos e desenvolvimento colaborativo de projetos em tempo real. Nesse contexto, as atividades que exploram essas tecnologias incluem:
\begin{itemize}
    \item Projetos integradores em equipes: uso de \textit{GitHub}, \textit{Google Colab} e \textit{Jupyter} para análise de dados e desenvolvimento de modelos em IA, com divisão de tarefas e colaboração;
    \item Aulas em tempo real: quizzes interativos, exercícios em grupo e discussões de problemas abertos via Moodle e Google Workspace;
Seminários e apresentações colaborativas: apresentações de trabalhos e projetos em videoconferência, integrando competências técnicas e comunicação;
    \item Laboratórios configurados para trabalho conjunto: mesas largas com múltiplos computadores para execução de experimentos em grupos; e,
    \item Fóruns assíncronos no Moodle: debate de conteúdos, troca de insights e construção coletiva de conhecimento.
\end{itemize}

A infraestrutura tecnológica é avaliada regularmente por meio da Comissão Própria de Avaliação (CPA), relatórios técnicos da COGETI-LD e comissões de curso, permitindo monitoramento da estabilidade, desempenho, adequação e qualidade de suporte aos usuários. Os resultados dessas avaliações geram planos documentados de atualização e aprimoramento, incluindo ajustes de hardware e software em laboratórios, atualização de plataformas digitais, capacitação de docentes e técnicos, implementação de novas funcionalidades e aquisição de ferramentas complementares. Este ciclo certifica que os recursos digitais atendam às necessidades do curso e proporcionem aprendizado seguro, atualizado e alinhado ao perfil profissional do egresso.

A COGETI-LD fornece suporte técnico a todas as plataformas digitais, incluindo Moodle, \textit{Google Workspace}, \textit{Google Colab}, \textit{GitHub Education}, Sistema Acadêmico e softwares de laboratório. O atendimento contempla resolução de problemas de login, configuração de softwares, manutenção de computadores, suporte à rede e atualização de ferramentas. Complementarmente, estudantes e docentes têm acesso à tutoria e capacitação pedagógica, por meio de cursos e workshops durante semanas acadêmicas, cursos internos assíncronos pelo AVA (Moodle), programas do Capacita Gov.br e Escola Virtual do Governo, além de monitorias em disciplinas com alta demanda tecnológica.

Para apoiar o aprendizado, o curso disponibiliza material didático digital interativo, incluindo tutoriais e exercícios online no Moodle, quizzes e avaliações interativas com feedback imediato, notebooks colaborativos via Google Colab e Jupyter, e integração com GitHub Education para versionamento e colaboração em projetos. Esses recursos fortalecem o aprendizado individual e colaborativo, promovendo competências técnicas, críticas e profissionais, alinhadas ao perfil do egresso.

\section{Acervo Bibliográfico}

O acervo bibliográfico é composto por materiais físicos e digitais, que atende às necessidades de ensino, pesquisa e extensão. A biblioteca do campus Londrina concede referências bibliográficas relacionadas às áreas de Ciência de Dados, Inteligência Artificial, Estatística, Matemática, Computação e áreas correlatas, disponíveis ao público acadêmico. O corpo técnico, formado por bibliotecários e assistentes administrativos, presta orientação na busca e recuperação de informações, tanto no acervo físico quanto em bases digitais, de modo que estudantes e docentes possam utilizar os recursos de forma eficiente.

Para manter a relevância do acervo físico, sua atualização ocorre por meio de um planejamento entre docentes e o colegiado do curso, que revisam os Planos de Ensino e solicitam novas edições ou títulos substitutos quando observada uma obra desatualizada. Dessa forma, a bibliografia básica e complementar indicada no PPC está disponível em quantidade adequada. O acesso ao acervo físico pode ser realizado durante todo o horário de funcionamento da biblioteca, de segunda a sexta-feira, das 08:00 às 21:30, permitindo tanto a consulta local quanto o empréstimo domiciliar de exemplares.

De acordo com o instrumento de avaliação do MEC, a definição das demandas de bibliografia básica e complementar é atribuição do Núcleo Docente Estruturante (NDE) do curso, realizada por meio do Plano de Adequação. Nesse contexto, as bibliotecas, diante das indicações do NDE, promovem aquisições pertinentes ao desenvolvimento das bibliografias básicas e complementares, nos formatos físico e digital, conforme previsto no relatório de adequação do acervo da biblioteca do campus Londrina.

A Política de Desenvolvimento de Coleções da UTFPR estabelece que o acervo deve crescer de forma planejada e estratégica, assegurando cobertura adequada das disciplinas ofertadas e atendimento às necessidades de acesso simultâneo e à projeção discente \cite{utfpr_ppdc}. Considerando que, no momento, não há acervo disponível específico para o novo curso, além dos registros no sistema MinhaBiblioteca  torna-se necessário que, com o início das atividades, o NDE elabore as demandas de acervo bibliográfico. A biblioteca poderá então realizar as aquisições necessárias, garantindo a conformidade com a política institucional e o crescimento planejado do acervo para suporte acadêmico adequado.

Complementando o acervo físico, o ecossistema digital da UTFPR oferece acesso integral a bases de dados científicas, periódicos eletrônicos e e-books por meio da Biblioteca Digital da UTFPR\footnote{https://www.utfpr.edu.br/bibliotecas/}. Estudantes e docentes têm acesso ao Portal de Periódicos CAPES, que inclui bases como IEEE Xplore, ACM Digital Library, ScienceDirect, SpringerLink e Scopus, bem como ao Repositório Institucional da UTFPR (RIUT), que reúne teses, dissertações e produção científica da universidade. Além disso, plataformas de e-books como Minha Biblioteca e Pearson disponibilizam milhares de títulos em português e inglês, abrangendo Estatística, Programação Python e R, Engenharia de Software, Banco de Dados e Aprendizado de Máquina, com leitura integral via navegador ou aplicativo e recursos de anotações e busca interna.

É importante destacar que diversas referências básicas das disciplinas do curso estão disponíveis no acervo digital proporcionando escalabilidade e flexibilidade de acesso aos alunos. O acesso remoto a esses recursos é viabilizado pela rede CAFe, permitindo que docentes e estudantes utilizem suas credenciais institucionais para acessar remotamente o acervo digital. O sistema de gestão Pergamum\footnote{https://biblioteca.utfpr.edu.br/} integra de forma híbrida os acervos físico e digital, permitindo que pesquisas sobre um tema específico apresentem simultaneamente livros, artigos e links diretos para e-books, promovendo maior eficiência na recuperação de informações.

As bases de periódicos internacionais disponibilizam artigos após publicação em congressos e revistas de referência, enquanto as plataformas de e-books incorporam novas edições e lançamentos de obras técnicas. Para apoiar o uso desses recursos, os estudantes recebem capacitação contínua, incluindo treinamentos em pesquisa avançada em bases de dados, utilização das bibliotecas digitais e aplicação das normas técnicas da ABNT. O serviço de referência, conduzido por bibliotecários qualificados, oferece auxílio na recuperação de informações, comutação bibliográfica e orientação para pesquisas em bases digitais.

O acervo está integrado às disciplinas, projetos de pesquisa e trabalhos de conclusão de curso, servindo como fonte para a fundamentação teórica de projetos de pesquisa, extensão e TCCs. Os docentes recomendam materiais de forma contínua, disponibilizando links diretos para e-books ou promovendo estudos dirigidos que exigem consulta a exemplares do acervo físico. No quesito acessibilidade, a biblioteca oferece recursos acessíveis para estudantes com deficiências visuais ou dificuldades de leitura, incluindo leitura em voz alta e ajustes de visualização em plataformas digitais. A biblioteca física também apresenta acessibilidade arquitetônica com rampas, corredores amplos, mobiliário adaptado e atendimento preferencial a portadores de necessidades especiais.

\section{Processo de Controle de Produção ou Distribuição de Material Didático (Logística)}

O curso segue um processo estruturado para a produção, distribuição, atualização e acompanhamento do material didático, promovendo qualidade pedagógica, acessibilidade e integração com o ensino semipresencial. A responsabilidade pela elaboração dos materiais é do professor da disciplina, que produz apostilas, grava videoaulas, disponibiliza artigos e tutoriais, propõe exercícios e compartilha arquivos por meio de plataformas como o \textit{Google Workspace} e Moodle.

A UTFPR orienta a produção dos materiais por meio do Manual de Identidade Visual \cite{utfpr2016manual} e das normas da ABNT para documentos científicos e técnicos, promovendo rigor acadêmico, uniformidade e padronização. No que diz respeito à acessibilidade, o Núcleo de Apoio às Pessoas com Necessidades Educacionais Específicas (NAPNE) atua como referência institucional, oferecendo aos estudantes com deficiência ou necessidades educacionais especiais, pleno acesso aos conteúdos. O NAPNE realiza adaptações de materiais complexos, oferece capacitação docente para a produção de documentos acessíveis e disponibiliza suporte contínuo aos estudantes.

O material didático será distribuído em formato digital, por meio do Ambiente Virtual de Aprendizagem institucional e de outras plataformas acadêmicas oficiais, com o objetivo de promover o acesso equitativo a todos os estudantes conforme preconiza a ODS 4 - Educação de Qualidade. A disponibilização digital permite que os conteúdos estejam acessíveis de forma contínua, organizada por unidades ou módulos, e estruturados conforme o cronograma estabelecido em cada disciplina. A adoção desse modelo também permite o controle e acompanhamento institucional dos acessos, possibilitando o monitoramento da utilização dos materiais, a verificação da participação discente e a geração de relatórios acadêmicos que subsidiam ações pedagógicas e administrativas. Além disso, a liberação gradual dos conteúdos, em consonância com o planejamento didático, favorece a organização do processo de aprendizagem, evitando sobrecarga informacional e promovendo maior engajamento dos estudantes.

O formato digital possibilita ainda a utilização de diferentes mídias como videoaulas, textos interativos, simuladores, estudos de caso, bases de dados e recursos multimídia, ampliando as estratégias pedagógicas e favorecendo metodologias ativas de ensino. Essa diversidade de recursos contribui para maior dinamismo, atualização constante do conteúdo e alinhamento às demandas contemporâneas da formação acadêmica. Adicionalmente, a priorização de materiais digitais busca minimizar a produção e a circulação de materiais impressos, contribuindo para a redução do consumo de papel atendendo a ODS 12 - Consumo e Produção Responsáveis \cite{UNICEF_ODS_2030}.

Os materiais são revisados e atualizados periodicamente, considerando os avanços do conteúdo de estudo. Os docentes também orientam continuamente os estudantes sobre o uso das plataformas digitais e os recursos disponíveis, especialmente no início de cada disciplina, para que se familiarizem com o ambiente virtual. A acessibilidade é promovida por meio de práticas específicas, como documentos estruturados, PDFs acessíveis, legendas e audiodescrição em videoaulas, além de recursos nativos das bibliotecas digitais, como leitura em voz alta e ajustes de contraste. O Moodle conta ainda com ferramentas que permitem adaptar a interface, ampliar fontes e integrar VLibras para tradução automática em Libras.

O material didático está integrado aos Planos de Ensino e às atividades avaliativas, sendo organizado por unidade, com referências bibliográficas, roteiros de laboratório, videoaulas e exercícios correspondentes. Essa integração oferece ao estudante um roteiro claro e completo para acompanhar a disciplina. O impacto do material no desempenho dos alunos é acompanhado por meio de métricas de engajamento no Moodle, incluindo relatórios de acesso, logs de atividade e barras de progresso, o que permite identificar conteúdos que exigem maior atenção. Além disso, o feedback direto obtido pela Comissão Própria de Avaliação (CPA) contribui para o aprimoramento contínuo da produção e distribuição do material didático.

\section{Laboratório Didático}

Os laboratórios didáticos do curso são estruturados para suportar atividades práticas das disciplinas de base e intermediárias, assegurando que os conteúdos teóricos sejam consolidados por meio de experimentação e aplicação prática. Os laboratórios disponíveis, com capacidade e equipamentos adequados ao número de estudantes e ao tipo de atividade são apresentados no Quadro \ref{quadro:4.3}.

\begin{quadro}[htb]
    \centering
    \caption{Características dos laboratórios didáticos}
    \label{quadro:4.3}
    \vspace{0.2cm}
    \begin{tabular}{|l|c|c|c|}
        \hline
        \textbf{Laboratório} & \textbf{Local} & \textbf{Área ($m^2$)} & \textbf{Capacidade (alunos)} \\ \hline
        Laboratório de Informática 1 & K-106 & 78 & 40 \\ \hline
        Laboratório de Informática 2 & K-107 & 78 & 40 \\ \hline
        Laboratório de Informática 3 & K-109 & 78 & 40 \\ \hline
        Laboratório de Informática 4 & K-102 & 25 & 24 \\ \hline
        Laboratório de Informática 5 & T-106 & 78 & 40 \\ \hline
        Laboratório de CAD & K-113 & 78 & 40 \\ \hline
        Laboratório de Eletrotécnica & K-007 & 80 & 24 \\ \hline
    \end{tabular}
    %\fonte{Autoria prória (2026)}
    \small \\ Fonte: Autoria própria (2026)
\end{quadro}

Nos laboratórios, os estudantes desenvolvem atividades como:

\begin{itemize}
    \item Programação em Python e R;
    \item Análise estatística de grandes volumes de dados;
    \item Construção e validação de algoritmos de aprendizado de máquina;
    \item Execução de exercícios de integração entre software e hardware em IA aplicada;
    \item Desenvolvimento de projetos práticos em banco de dados, visualização de dados e prototipagem de sistemas inteligentes.
\end{itemize}

A gestão dos laboratórios inclui planejamento de horários, supervisão docente, manutenção periódica, atualização de softwares e controle da ocupação, garantindo que cada espaço esteja plenamente adequado para o desenvolvimento das atividades previstas no currículo. O foco deste item é evidenciar estrutura, equipamentos e condições de uso, demonstrando que os laboratórios estão preparados para apoiar integralmente o aprendizado prático dos estudantes.

\section{Biblioteca}

A Biblioteca Central da UTFPR \textit{campus} Londrina é integrante do Sistema de Bibliotecas da instituição, composto por 15 unidades distribuídas nos 13 campi. A biblioteca atua de forma estratégica no apoio ao ensino, à pesquisa, à extensão e à inovação, disponibilizando produtos e serviços informacionais alinhados às demandas acadêmicas e científicas contemporâneas.

A unidade que atende ao curso possui área física aproximada de 415 m², organizada para proporcionar ambientes adequados ao estudo individual e coletivo. Dispõe de salas para estudo individual e em grupo, equipadas com mobiliário apropriado, acesso à internet cabeada e rede sem fio (Wi-Fi), além de espaços de convivência acadêmica que favorecem a permanência e o desenvolvimento das atividades discentes. O horário de funcionamento é de segunda a sexta-feira, das 8h às 21h30, garantindo atendimento contínuo nos turnos diurno e noturno.

O acervo necessário ao curso encontra-se disponível em formato digital, assegurando atendimento integral à bibliografia básica e complementar prevista no PPC, com no mínimo três títulos por unidade curricular, conforme diretrizes institucionais. A política de desenvolvimento e atualização do acervo é formalizada e ocorre de maneira contínua, especialmente por meio de parcerias com plataformas consolidadas, tais como Minha Biblioteca, EBSCO e outras bases especializadas, garantindo atualização frequente e aderência às evoluções científicas nas áreas de Ciência de Dados, Inteligência Artificial, Estatística, Computação, Matemática e áreas correlatas.

No que se refere ao acervo digital e às bases de dados científicas, a instituição disponibiliza:

\begin{itemize}
    \item Biblioteca virtual (Minha Biblioteca, EBSCO, entre outras);
    \item Acesso integral ao Portal de Periódicos CAPES;
    \item Bases de dados de referência internacional, como IEEE, ACM Digital Library, Scopus, Web of Science, além de editoras científicas como Elsevier, Springer e MDPI.
\end{itemize}

O acesso remoto é realizado por meio da Comunidade Acadêmica Federada (CAFe), possibilitando que docentes e discentes consultem livros, periódicos e bases especializadas fora do ambiente físico da instituição, ampliando as condições de estudo, pesquisa e produção científica.

A biblioteca oferece suporte qualificado por meio de equipe técnico-administrativa composta por bibliotecária e técnicos especializados, garantindo atendimento contínuo à comunidade acadêmica. São ofertadas capacitações presenciais e a distância sobre uso de bases científicas, normalização de trabalhos acadêmicos conforme normas da ABNT, e gestão de referências bibliográficas com ferramentas como Mendeley e Zotero. Também são disponibilizados serviços de levantamento informacional, apoio à revisão sistemática, comutação bibliográfica, elaboração de ficha catalográfica e disseminação seletiva da informação.

O gerenciamento do acervo físico é realizado por meio do sistema informatizado Pergamum, que permite consulta ao catálogo, empréstimos, renovações e emissão de relatórios gerenciais de uso, constituindo importante instrumento de acompanhamento e planejamento. O metabuscador institucional (Bibliotec) possibilita consulta integrada aos acervos físico e digital, bem como a realização de empréstimos entre bibliotecas da UTFPR e de outras instituições.

No âmbito da acessibilidade, a biblioteca dispõe de infraestrutura física adequada, incluindo rampas de acesso, espaços adaptados para cadeirantes e organização do mobiliário que assegura circulação adequada, reafirmando o compromisso institucional com a inclusão e a permanência estudantil.

A biblioteca integra os processos de avaliação institucional conduzidos pela Comissão Própria de Avaliação (CPA) \cite{UTFPR_CPA_2004, UTFPR_CPA_2009}, utilizando indicadores de uso e satisfação para aprimoramento contínuo dos serviços ofertados. Dessa forma, a Biblioteca da UTFPR configura-se como infraestrutura plenamente adequada às necessidades do curso de CDIA, concedendo suporte informacional atualizado, acesso a bases científicas internacionais, serviços especializados de apoio acadêmico e condições estruturais compatíveis com a formação técnico-científica exigida na área.